% Project Milestone Report — CVPR 2022 Template
\documentclass[10pt,twocolumn,letterpaper]{article}

\usepackage[pagenumbers]{cvpr}

\usepackage{hyperref}
\usepackage{multirow}

\def\paperID{N/A}
\def\confName{CVPR}
\def\confYear{2022}

\title{Cross-Feature Attention for Architecture-Level CPU Power Modeling:\\An FT-Transformer Approach on the ArchPower Dataset}

\author{Jingbo Jiang}

\begin{document}
\maketitle

%% ====================================================================
\begin{abstract}
Accurate architecture-level power estimation is critical for early-stage
CPU design-space exploration.  We propose using the Feature Tokenizer
Transformer (FT-Transformer) for component-level CPU power prediction.
FT-Transformer tokenizes each of 81 event features into a learned
embedding and applies multi-head self-attention to model
\emph{cross-feature interactions}---capturing how coupled
microarchitectural activities jointly determine power.
On the ArchPower dataset (200 samples, 2 RISC-V architectures),
FT-Transformer achieves 3.61\%/5.33\% MAPE (BOOM/XiangShan), outperforming
Ridge (4.36\%/5.95\%) and MLP (6.29\%/7.76\%) baselines.  While XGBoost
currently edges ahead (2.36\%/4.26\%) due to its small-data advantage,
we propose collecting an expanded dataset for the final report where we
expect FT-Transformer to achieve the best accuracy.
\end{abstract}

%% ====================================================================
\section{Introduction}

Power consumption is the primary design constraint for modern processors.
Accurate power evaluation traditionally requires the full RTL
implementation, logic synthesis, and signoff-grade power analysis—a process
that easily takes months and precludes rapid design-space exploration at
the architecture level~\cite{zhang2025archpower}.  Classical analytical
tools such as McPAT~\cite{li2009mcpat} attempt to fill this gap but are
known to exhibit large modeling errors on modern out-of-order CPUs.

Recent work has proposed \emph{ML-calibrated} analytical models.
McPAT-Calib~\cite{zhai2021mcpat_calib} trains XGBoost regressors to correct
McPAT's component-level estimates, while PANDA~\cite{zhang2023panda}
introduces physics-aware \emph{resource functions} that encode
hardware-dependent scaling prior to ML training.  Neural-network-based
approaches have also been explored~\cite{zhai2023neural_power}, but
existing methods treat input features independently—missing the rich
\emph{cross-feature interactions} inherent in CPU power behavior.  For
example, branch misprediction rates interact with instruction fetch
bandwidth to determine front-end power, and cache miss rates couple with
memory controller traffic to drive back-end power.  Models that capture
these pairwise interactions should yield both better predictions and more
interpretable insights.

We propose applying the Feature Tokenizer Transformer
(FT-Transformer)~\cite{gorishniy2021ft_transformer} to architecture-level
power modeling.  FT-Transformer tokenizes each scalar feature into a
learned embedding, then applies multi-head self-attention over all feature
tokens.  This design is uniquely suited to power modeling because:
(1)~each attention head learns a distinct interaction pattern (\eg,
pipeline-stage coupling \vs memory-hierarchy coupling);
(2)~the $81 \times 81$ attention matrices are directly interpretable as
event-to-event interaction maps; and
(3)~multi-target regression provides implicit regularization critical for
the small-data regime.

Our contributions at this milestone are:
(1)~an \textbf{event-only preprocessing} pipeline that factors out
hardware dependence via PANDA's resource functions (101$\to$81 features);
(2)~demonstrating that \textbf{FT-Transformer outperforms} both Ridge and
MLP baselines, validating cross-feature attention for power prediction; and
(3)~a \textbf{data scaling hypothesis} arguing that XGBoost's current lead
is an artifact of the limited dataset size.

%% ====================================================================
\section{Problem Statement}

We use the ArchPower dataset~\cite{zhang2025archpower}, an open-source
benchmark with 200 CPU data samples from two RISC-V architectures:
\textbf{BOOM} (15 configs $\times$ 8 workloads = 120 samples) and
\textbf{XiangShan} (10 configs $\times$ 8 workloads = 80 samples).
Each sample has a 101-d feature vector (20 hardware parameters + 81 event
parameters) and power labels for 12 components $\times$ 5 power groups.

\noindent\textbf{Task.}  Given only the 81 event features (after removing
hardware dependence), predict the 12-component power vector for unseen CPU
configurations.  Splits are at the \emph{configuration level}---all 8
workloads from a config are entirely in train or test---to evaluate
cross-configuration generalization.

\noindent\textbf{Evaluation.}  MAPE, $R^2$, and RMSE on a config-level
holdout (20\% of configs).

\noindent\textbf{Challenge.}  The extreme sample-to-feature ratio (120 or
80 samples for 81 features) makes overfitting the dominant risk.

%% ====================================================================
\section{Technical Approach}

\subsection{Data Preprocessing: Event-Only Dataset}
\label{sec:preproc}
Following PANDA~\cite{zhang2023panda}, we model each component's power as
$P_c = R_c(\mathbf{x}_\text{hw}) \cdot f_c(\mathbf{x}_\text{event})$ and
divide out the hardware-dependent resource function $R_c$ to obtain
normalized labels $P'_c = P_c / R_c$.  Resource functions are
\emph{multiplicative} for most components (BP, ICache, IFU, RNU, LSU,
DCache, ROB), \emph{additive} for Regfile, \emph{lookup-table} for ISU,
and \emph{identity} for FU-Pool.
After normalization, the 20 hardware-constant columns are removed, yielding
an 81-d event-only feature matrix used by all downstream models.

\subsection{Proposed Model: FT-Transformer}

The core of our approach is the Feature Tokenizer Transformer
(FT-Transformer)~\cite{gorishniy2021ft_transformer}, designed specifically
for tabular data.  Unlike MLPs that process the concatenated feature vector
through dense layers, FT-Transformer \emph{tokenizes each feature
independently}: every scalar event feature $x_j$ is projected into a
$d$-dimensional embedding $\mathbf{t}_j = \mathbf{w}_j x_j +
\mathbf{b}_j$.  A \texttt{[CLS]} token is appended, and the 82 tokens
are processed by transformer blocks with multi-head self-attention.

\noindent\textbf{Why attention captures cross-feature interactions.}
Power consumption arises from complex interactions between
microarchitectural activities---\eg, instruction cache misses increase
fetch stalls, reducing issue-queue utilization and functional-unit power.
Linear models ignore these interactions; MLPs learn them only implicitly.
In contrast, each self-attention head computes an explicit
$81\times 81$ interaction matrix
$\mathbf{A} = \text{softmax}(\mathbf{Q}\mathbf{K}^\top / \sqrt{d_k})$,
where $A_{ij}$ quantifies how strongly event $i$ attends to event $j$.
Multiple heads capture distinct patterns simultaneously---pipeline-stage
coupling, memory-hierarchy dependencies, etc.

\noindent\textbf{Architecture and training.}
We use $d_\text{token}{=}32$ (BOOM) or 48 (XS), 2 transformer blocks,
4 attention heads, and FFN dimension 64 (22K--40K parameters).
Training has two phases: (1)~\emph{self-supervised pre-training} via masked
feature modeling, where random features are masked and reconstructed to
learn feature correlations; (2)~\emph{supervised fine-tuning} with MSE loss
on all 12 power components and early stopping.  After training, we extract
attention weight matrices to construct event-to-event interaction maps.

\subsection{Baseline Models}

\noindent\textbf{Ridge Regression.}  Multi-output Ridge
(81$\to$12, 984 params)---captures only first-order relationships.

\noindent\textbf{3-Layer MLP.}  Hidden layers (128, 64, 32) with ReLU and
L2 regularization (21K params)---learns nonlinear mappings but models
interactions only implicitly.

\noindent\textbf{XGBoost GBDT.}  Twelve independent regressors (100 trees,
depth 6)---excels on small tabular data but lacks explicit cross-feature
interaction modeling and interpretability.

%% ====================================================================
\section{Intermediate/Preliminary Results}

\subsection{FT-Transformer Outperforms Ridge and MLP}

Table~\ref{tab:results} summarizes total-power prediction accuracy.  Our
FT-Transformer substantially outperforms both baselines: it reduces MAPE
by 17\% relative to Ridge (3.61\% \vs 4.36\% on BOOM) and by 43\% relative
to the MLP (3.61\% \vs 6.29\%).  FT-Transformer achieves $R^2$ of 0.925
(BOOM) and 0.932 (XS), compared to 0.856/0.905 for Ridge and 0.629/0.858
for the MLP.  These gains validate that \emph{explicit cross-feature
attention improves power prediction} over models that either ignore
interactions (Ridge) or learn them only implicitly (MLP).

\begin{table}[t]
\centering
\small
\caption{Total-power prediction on BOOM and XS (config-level holdout,
val\_ratio=0.2).  Best deep learning result \underline{underlined};
best overall in \textbf{bold}.}
\label{tab:results}
\begin{tabular}{l cc cc}
\toprule
& \multicolumn{2}{c}{\textbf{BOOM}} & \multicolumn{2}{c}{\textbf{XS}} \\
\cmidrule(lr){2-3} \cmidrule(lr){4-5}
\textbf{Model} & MAPE\% & $R^2$ & MAPE\% & $R^2$ \\
\midrule
Ridge        & 4.36 & 0.856 & 5.95 & 0.905 \\
MLP          & 6.29 & 0.629 & 7.76 & 0.858 \\
\midrule
FT-Trans.    & \underline{3.61} & \underline{0.925} & \underline{5.33} & \underline{0.932} \\
XGBoost      & \textbf{2.36} & \textbf{0.938} & \textbf{4.26} & \textbf{0.945} \\
\bottomrule
\end{tabular}
\end{table}

\subsection{Component-Level Analysis}

Table~\ref{tab:component} compares per-component MAPE on BOOM.  While
XGBoost leads on most components, FT-Transformer \emph{outperforms XGBoost
on three components}: ROB (18.25\% \vs 25.75\%), RNU (28.71\% \vs
39.18\%), and Others (30.51\% \vs 42.94\%).  These components' power
depends on complex inter-unit interactions (\eg, ROB occupancy couples
with rename width and issue-queue pressure), suggesting that attention
provides a genuine advantage where cross-feature interactions dominate.

\begin{table}[t]
\centering
\small
\caption{Selected per-component MAPE (\%) on BOOM.  Components where
FT-Transformer outperforms XGBoost in \textbf{bold}.}
\label{tab:component}
\begin{tabular}{l rr}
\toprule
\textbf{Component} & \textbf{XGBoost} & \textbf{FT-Trans.} \\
\midrule
Total    &  2.36 &  3.61 \\
BP       &  1.05 &  1.04 \\
Regfile  & 13.24 & 14.34 \\
ROB      & 25.75 & \textbf{18.25} \\
RNU      & 39.18 & \textbf{28.71} \\
Others   & 42.94 & \textbf{30.51} \\
\bottomrule
\end{tabular}
\end{table}

\subsection{Why XGBoost Leads---and Why This Will Change}

XGBoost currently achieves the best total-power MAPE (2.36\% / 4.26\%),
outperforming FT-Transformer by $\sim$1.3 percentage points.  We attribute
this to a well-known phenomenon: \emph{tree-based ensembles dominate
neural models in the small-data tabular
regime}~\cite{gorishniy2021ft_transformer}.  With only 200 samples,
gradient-boosted trees benefit from built-in feature selection and
sample-efficient splitting.  The original FT-Transformer
paper~\cite{gorishniy2021ft_transformer} shows that this gap
\emph{narrows and reverses} as dataset size grows, motivating our plan to
collect a larger dataset for the final report.

\subsection{Next Steps}

For the final report, we plan to:
(1)~\textbf{Collect an expanded dataset} with additional CPU configurations
and workloads (target 1{,}000+ samples), testing whether FT-Transformer
surpasses XGBoost as expected;
(2)~\textbf{Analyze attention maps} to identify previously unknown event
interactions for CPU designers; and
(3)~\textbf{Refine the model} with deeper variants and architectures that
enforce the component-to-total summation constraint.

{\small
\bibliographystyle{ieeenat_fullname}
\bibliography{references}
}

\end{document}
